
\bibitem{raymond1999}
Eric S. Raymond. 1999. \emph{The Cathedral and the Bazaar: Musings on Linux and Open Source by an Accidental Revolutionary}. O'Reilly.

\bibitem{omahony2007}
Siobh\'an O'Mahony and Fabrizio Ferraro. 2007. The emergence of governance in an open source community. \emph{Academy of Management Journal} 50, 5, 1079--1106.

\bibitem{dahlander2005}
Linus Dahlander and Mats G. Magnusson. 2005. Relationships between open source software companies and communities: Observations from Nordic firms. \emph{Research Policy} 34, 4, 481--493.

\bibitem{riehle2007}
Dirk Riehle. 2007. The economic motivation of open source software: Stakeholder perspectives. \emph{IEEE Computer} 40, 4, 25--32.

\bibitem{zhang2020}
Yuxia Zhang, Minghui Zhou, Klaas-Jan Stol, Jianyu Wu, and Zhi Jin. 2020. How do companies collaborate in open source ecosystems? An empirical study of OpenStack. In \emph{Proceedings of the 42nd International Conference on Software Engineering (ICSE 2020)}, 1196--1208.

\bibitem{ducheneaut2005}
Nicolas Ducheneaut. 2005. Socialization in an open source software community: A socio-technical analysis. \emph{Computer Supported Cooperative Work} 14, 4, 323--368.

\bibitem{jensen2007}
Chris Jensen and Walt Scacchi. 2007. Role migration and advancement processes in OSSD projects: A comparative case study. In \emph{Proceedings of the 29th International Conference on Software Engineering (ICSE 2007)}, 364--374.

\bibitem{hirschman1970}
Albert O. Hirschman. 1970. \emph{Exit, Voice, and Loyalty: Responses to Decline in Firms, Organizations, and States}. Harvard University Press.

\bibitem{robles2012}
Gregorio Robles and Jes\'us M. Gonz\'alez-Barahona. 2012. A comprehensive study of software forks: Dates, reasons and outcomes. In \emph{Open Source Systems: Long-Term Sustainability (OSS 2012)}, IFIP AICT 378, 1--14.

\bibitem{gamalielsson2014}
Jonas Gamalielsson and Bj\"orn Lundell. 2014. Sustainability of open source software communities beyond a fork: How and why has the LibreOffice project evolved? \emph{Journal of Systems and Software} 89, 128--145.

\bibitem{raman2020}
Naveen Raman, Minxuan Cao, Yulia Tsvetkov, Christian K\"astner, and Bogdan Vasilescu. 2020. Stress and burnout in open source: Toward finding, understanding, and mitigating unhealthy interactions. In \emph{ICSE-NIER 2020}, 57--60.

\bibitem{ferreira2021}
Isabella Ferreira, Jinghui Cheng, Bram Adams, and Eirini Kalliamvakou. 2021. The ``shut the f**k up'' phenomenon: Characterizing incivility in open source code review discussions. \emph{Proceedings of the ACM on Human-Computer Interaction} 5, CSCW2, Article 353.

\bibitem{squire2015}
Megan Squire and Rebecca Gazda. 2015. FLOSS as a source for profanity and insults: Collecting the data. In \emph{Proceedings of the 48th Hawaii International Conference on System Sciences (HICSS 2015)}, 5290--5298.

\bibitem{schneider2016}
Daniel Schneider, Scott Spurlock, and Megan Squire. 2016. Differentiating communication styles of leaders on the Linux kernel mailing list. In \emph{Proceedings of the 12th International Symposium on Open Collaboration (OpenSym 2016)}, Article 2.

\bibitem{delaat2007}
Paul B. de Laat. 2007. Governance of open source software: State of the art. \emph{Journal of Management \& Governance} 11, 2, 165--177.

\bibitem{fitzgerald2006}
Brian Fitzgerald. 2006. The transformation of open source software. \emph{MIS Quarterly} 30, 3, 587--598.

\bibitem{transferofpower}
Guido van Rossum. 2018. Transfer of power. Message to python-committers, 12 July 2018. Retrieved September 12, 2026 from \url{https://mail.python.org/pipermail/python-committers/2018-July/005664.html}
